% ============================================================
%  real_time_optimization_model.tex
%  Multi-Tenant Cluster Scheduling — Mathematical Formulation
%  DAMO 699 Capstone · Group 3
% ============================================================
\documentclass[11pt]{article}

\usepackage[margin=1in]{geometry}
\usepackage{amsmath}
\usepackage{amssymb}
\usepackage{mathtools}
\usepackage{booktabs}
\usepackage{array}
\usepackage{enumitem}
\usepackage[T1]{fontenc}
\usepackage{lmodern}
\usepackage{microtype}
\usepackage[colorlinks=true,linkcolor=blue,urlcolor=blue]{hyperref}

\setlength{\parskip}{0.5em}
\setlength{\parindent}{0pt}

% ── convenience macros ────────────────────────────────────────────────────────
\newcommand{\vbarsla}{\bar{v}_n^{\text{SLA}}}
\newcommand{\Meff}{M_n^{\text{eff}}}
\newcommand{\Mavail}{M_n^{\text{avail}}}
\newcommand{\Mtax}{M_n^{\text{tax}}}
\newcommand{\Mremain}{M_n^{\text{remain}}}
\newcommand{\omfixed}{\omega_n^{\text{fixed}}}
\newcommand{\umem}{u_n^{\text{mem}}}
\newcommand{\Pmem}{\hat{P}_j^{\text{mem}}}
\newcommand{\Pcpu}{\hat{P}_j^{\text{CPU}}}
\newcommand{\Plife}{\hat{P}_j^{\text{life}}}

% ── title ─────────────────────────────────────────────────────────────────────
\title{\textbf{Real-Time Optimization Model}\\[0.5em]
       \large Mathematical Formulation}
\author{DAMO 699 Capstone $\cdot$ Group 3\\
        Multi-Tenant Cluster Memory Scheduling}
\date{}

\begin{document}
\maketitle
\tableofcontents
\newpage

% =============================================================================
\section*{Terminology}
\addcontentsline{toc}{section}{Terminology}
% =============================================================================

A \textbf{job} refers to any containerized workload unit submitted by a tenant
for execution on the cluster. In Kubernetes this corresponds to a \textbf{pod}
--- the smallest schedulable unit. We use \emph{job} for generality, as the
model is scheduler-agnostic and not specific to Kubernetes.

% =============================================================================
\section{System Architecture}
% =============================================================================

The cluster monitor provides real-time inputs to the optimizer each scheduling
round and receives the placement assignment $x_{jn}$ in return.

\medskip
\noindent\textbf{Cluster monitor inputs:}
\begin{itemize}[noitemsep, topsep=2pt]
  \item $U_n$ --- currently used memory per node
  \item $\vbarsla$ --- rolling SLA violation rate per node
  \item $\omega_j$ --- priority weight per job (inherited from the job's tenant)
\end{itemize}

\medskip
\noindent\textbf{Prediction layer inputs:}
\begin{itemize}[noitemsep, topsep=2pt]
  \item $\Pmem$ --- P95 predicted memory per pending job
  \item $\Pcpu$ --- P90 predicted CPU peak per pending job
\end{itemize}

\medskip
The optimizer maximizes $Z$ (Placement Score) and outputs $x_{jn}$, which
the cluster monitor uses to submit jobs to nodes.

% =============================================================================
\section{Decision Variable}
% =============================================================================

\[
  x_{jn} \;\in\; \{0,1\}
\]

$x_{jn} = 1$ if job $j$ is assigned to node $n$; $x_{jn} = 0$ otherwise.
This placement matrix is the sole output of the optimizer.

% =============================================================================
\section{Sets}
% =============================================================================

\begin{center}
\begin{tabular}{@{}l p{10.5cm}@{}}
\toprule
Symbol & Definition \\
\midrule
$J$   & All jobs currently in the pending queue at this scheduling round \\[3pt]
$N$   & All nodes available in the cluster \\[3pt]
$T$   & All tenants sharing the cluster \\[3pt]
$J_t$ & Jobs belonging to tenant $t$, where $J_t \subseteq J$ and
        $\displaystyle\bigcup_{t \in T} J_t = J$ \\
\bottomrule
\end{tabular}
\end{center}

% =============================================================================
\section{Parameters}
% =============================================================================

% ── Node Specification ───────────────────────────────────────────────────────
\subsection*{Node Specification \normalfont\textit{[hardware]}}
\textit{Physical node properties --- fixed facts about the hardware.}

\medskip
\begin{center}
\begin{tabular}{@{}l p{10.5cm}@{}}
\toprule
Symbol & Meaning \\
\midrule
$M_n$ & Physical RAM installed on node $n$ \\[3pt]
$C_n$ & Total CPU cores on node $n$ \\
\bottomrule
\end{tabular}
\end{center}

% ── Configuration ────────────────────────────────────────────────────────────
\subsection*{Configuration \normalfont\textit{[config]}}
\textit{Set by the cluster operator before the simulation; constant at runtime.}

\medskip
\begin{center}
\begin{tabular}{@{}l p{10.5cm}@{}}
\toprule
Symbol & Meaning \\
\midrule
$\Mtax$    & \textbf{Memory tax} --- fixed overhead reserved for the OS,
  kubelet, and monitoring agents; not available to tenant jobs \\[6pt]
$\theta_n$ & \textbf{Memory safety threshold} --- reserves a fixed fraction
  of remaining capacity as a buffer against runtime memory spikes.
  $\theta_n \in [0,1]$. Default: $\theta_n = 0.10$ (10\% reserved) \\[6pt]
$K$        & \textbf{Rolling window length} --- number of past scheduling
  rounds used to compute $\vbarsla$. See Appendix~\ref{app:vbar} \\[6pt]
$\omfixed$ & \textbf{Fixed node consolidation weight} --- static bias that
  steers placement toward lower-indexed nodes, keeping the cluster
  consolidated. Node 0 gets the highest weight; the last node gets 1.
  Constant across all scheduling rounds.
  $\omfixed = |N| - n,\quad n = 0, 1, \ldots, |N|-1$

  \smallskip\textit{Note:} A sequencing constraint (node $n$ may only receive
  jobs if node $n{-}1$ is non-empty) would create a recursive dependency the
  MILP cannot resolve in a single solve call. The fixed weight achieves the
  same consolidation bias without constraint coupling. \\
\bottomrule
\end{tabular}
\end{center}

% ── Cluster Monitor Inputs ───────────────────────────────────────────────────
\subsection*{Cluster Monitor Inputs \normalfont\textit{[monitor]}}
\textit{Provided to the optimizer at the start of each scheduling round.}

\medskip
\begin{center}
\begin{tabular}{@{}l p{10.5cm}@{}}
\toprule
Symbol & Meaning \\
\midrule
$U_n$ & Currently used memory on node $n$, measured in real time \\[6pt]
$\vbarsla$ & \textbf{Rolling SLA violation rate} on node $n$ --- fraction of
  the last $K$ rounds in which actual memory on node $n$ exceeded $\Mavail$.
  $\bar{v}_n^{\text{SLA}} \in [0,1]$.
  $\vbarsla = 0$: no recent violations. \quad
  $\vbarsla = 1$: every recent round was overloaded.
  See Appendix~\ref{app:vbar} \\[6pt]
$\omega_j$ & \textbf{Priority weight} for job $j$ --- inherited from the
  tenant that owns job $j$. $\omega_j \geq 1$.

  \smallskip
  $\omega_j = 1$: no extra weight (tenant wait at or below cluster average).\\
  $\omega_j > 1$: additional weight; the solver increases preference for job
  $j$ (real-valued, no maximum). See Appendix~\ref{app:omega} \\
\bottomrule
\end{tabular}
\end{center}

% ── Prediction Layer Inputs ──────────────────────────────────────────────────
\subsection*{Prediction Layer Inputs \normalfont\textit{[prediction]}}
\textit{Provided by the prediction layer for each pending job.}

\medskip
\begin{center}
\begin{tabular}{@{}l p{10.5cm}@{}}
\toprule
Symbol & Meaning \\
\midrule
$\Pmem$  & P95 predicted memory for job $j$ --- 95th percentile of the
  highest predicted memory spike over the job's expected lifetime. Used in
  constraint C2 and the objective \\[6pt]
$\Pcpu$  & P90 predicted CPU peak for job $j$ --- 90th percentile of the
  job's predicted CPU demand. Used in constraint C4. No SLA feedback applies
  to CPU; kernel throttling handles runtime overrun \\[6pt]
$\Plife$ & Predicted job lifetime. Used by the prediction layer to determine
  the time horizon for $\Pmem$ and $\Pcpu$. Not used directly in the LP \\
\bottomrule
\end{tabular}
\end{center}

% ── Derived ──────────────────────────────────────────────────────────────────
\subsection*{Derived \normalfont\textit{[derived]}}
\textit{Computed by the optimizer from the above inputs each scheduling round.}

\medskip
\begin{center}
\begin{tabular}{@{}l p{10.5cm}@{}}
\toprule
Symbol & Meaning \\
\midrule
$\Mavail$ & Available memory --- physical RAM minus the memory tax; fixed
  upper bound for tenant job memory on node $n$ \\[6pt]
$\Mremain$ & Remaining available memory after subtracting current usage \\[6pt]
$\Meff$ & \textbf{Effective remaining memory} --- capacity offered to new
  jobs after applying the safety threshold and SLA violation scaling. This is
  the right-hand side of constraint C2 \\[6pt]
$\umem$ & \textbf{Memory utilization weight} --- fraction of $\Mavail$
  currently in use on node $n$, mapped to $[1,2]$. Applied in the objective
  to consolidate workloads onto busier nodes.
  Constraint C2 prevents infeasible placements regardless of $\umem$.
  See Appendix~\ref{app:umem} \\
\bottomrule
\end{tabular}
\end{center}

\begin{align}
  \Mavail  &\;=\; M_n - \Mtax \notag \\[4pt]
  \Mremain &\;=\; \Mavail - U_n \notag \\[4pt]
  \Meff    &\;=\; \max\!\left(0,\;\;
               \Mremain \cdot (1 - \theta_n) \cdot
               (1 - \bar{v}_n^{\text{SLA}})\right) \notag
\end{align}

% =============================================================================
\section{Feedback Loops}
% =============================================================================

SLA and fairness are not objectives inside the LP --- they are observed
outcomes tracked by the cluster monitor and fed back as adaptive inputs each
scheduling round.

\paragraph{SLA Feedback $\rightarrow$ Node Capacity.}
When actual memory on node $n$ exceeds $\Mavail$ at runtime, the cluster
monitor records one overload event. It computes $\vbarsla$ over the last $K$
rounds and passes it to the optimizer. As violations persist, $\vbarsla$ rises
and $\Meff$ shrinks --- the node accepts fewer new jobs until violations
subside. One overload event equals one SLA violation; the model reacts
conservatively but cannot retroactively prevent an overload.

\paragraph{Fairness Feedback $\rightarrow$ Job Priority.}
The cluster monitor tracks per-job wait times, computes $\omega_j$ per job,
and passes it to the optimizer. Jobs whose tenant is waiting above average
receive $\omega_j > 1$, so the solver naturally prefers placing them. Fairness
improves as a side effect of weighted maximization. Variance of $\bar{W}_t$
across tenants (or Gini coefficient) is reported as a simulation KPI --- not
minimized inside the LP.

% =============================================================================
\section{Optimization Objective}
% =============================================================================

Maximize the \textbf{Placement Score} $Z$:

\[
  \max \quad Z \;=\;
  \sum_{j \in J}\; \sum_{n \in N}
    \omega_j \cdot \Pmem \cdot \umem \cdot \omfixed \cdot x_{jn}
\]

\noindent\textit{Note on $\omega_j$:} $\omega_j = \omega_{t(j)}$, the
priority weight of the tenant that owns job $j$, where $t(j)$ denotes that
tenant. Writing $\omega_j$ rather than $\omega_{t(j)}$ is standard in
weighted scheduling formulations where the job weight is inherited from its
class or owner. It is mathematically equivalent.

\medskip
Each placed job contributes its P95 predicted memory $\Pmem$, scaled by three
multiplicative weights:
\begin{itemize}[noitemsep, topsep=4pt]
  \item $\omega_j \geq 1$ --- tenant priority: jobs from long-waiting tenants
    get higher weight, so the solver prefers placing them.
  \item $\umem \in [1,2]$ --- memory consolidation: a job placed on a
    memory-loaded node contributes more, consolidating workloads onto fewer
    machines.
  \item $\omfixed \in \{1,\ldots,|N|\}$ --- fixed node bias: node 0 has the
    highest weight, ensuring consolidation even at the first batch when all
    $u_n^{\text{mem}} = 1$ (no prior load history).
\end{itemize}

% =============================================================================
\section{Constraints}
% =============================================================================

\paragraph{C1 --- One node per job.}
Each job is assigned to at most one node per scheduling round:
\[
  \sum_{n \in N} x_{jn} \;\leq\; 1 \qquad \forall\, j \in J
\]
$\leq 1$ (not $= 1$): a job may remain unscheduled if no node has sufficient
capacity. Unscheduled jobs return to the queue, accumulating wait time that
raises their $\omega_j$ in subsequent rounds.

\paragraph{C2 --- Memory capacity per node.}
Total P95 predicted memory of jobs assigned to node $n$ may not exceed its
effective remaining capacity:
\[
  \sum_{j \in J} \Pmem \cdot x_{jn} \;\leq\; \Meff \qquad \forall\, n \in N
\]
This is a hard constraint. $\Meff$ incorporates the OS tax, current usage,
safety threshold, and SLA-driven scaling. The model admits a job if $\Pmem$
fits within $\Meff$ even if the tenant's declared request would not
(overcommitment). An SLA violation occurs only if the 5\% prediction tail
materializes at runtime.

\paragraph{C3 --- Binary domain.}
\[
  x_{jn} \in \{0,1\} \qquad \forall\, j \in J,\; n \in N
\]
This makes the problem a Mixed-Integer Linear Program (MILP). For large
instances, $x_{jn} \in [0,1]$ (LP relaxation with rounding) is an acceptable
approximation.

\paragraph{C4 --- Per-pair CPU fitment.}
\[
  x_{jn} = 0 \;\text{ if }\; \Pcpu > C_n
  \qquad \forall\, j \in J,\; n \in N
\]
A job may not be placed on a node whose total core count is less than the
job's predicted CPU peak. CPU cores are time-shared --- the kernel handles
overcommit via throttling. This check blocks only the genuinely infeasible
case where a job's peak CPU demand exceeds the node's hardware ceiling.

% =============================================================================
\section{Complete Model}
% =============================================================================

\[
  \max \quad Z \;=\;
  \sum_{j \in J}\; \sum_{n \in N}
    \omega_j \cdot \Pmem \cdot \umem \cdot \omfixed \cdot x_{jn}
\]

\noindent\textbf{Subject to:}
\begin{align*}
  \text{C1:}&\quad \sum_{n \in N} x_{jn} \leq 1
    && \forall\, j \in J \\[4pt]
  \text{C2:}&\quad \sum_{j \in J} \Pmem \cdot x_{jn} \leq \Meff
    && \forall\, n \in N \\[4pt]
  \text{C3:}&\quad x_{jn} \in \{0,1\}
    && \forall\, j \in J,\; n \in N \\[4pt]
  \text{C4:}&\quad x_{jn} = 0 \;\text{ if }\; \Pcpu > C_n
    && \forall\, j \in J,\; n \in N
\end{align*}

\noindent\textbf{Where:}
\begin{align*}
  \Mavail  &\;=\; M_n - \Mtax \\[4pt]
  \Mremain &\;=\; \Mavail - U_n \\[4pt]
  \Meff    &\;=\; \max\!\left(0,\;\;
               \Mremain \cdot (1 - \theta_n) \cdot
               (1 - \bar{v}_n^{\text{SLA}})\right) \\[4pt]
  \omfixed &\;=\; |N| - n, \quad n = 0,\,1,\,\ldots,\,|N|-1
\end{align*}

\medskip
\begin{center}
\small
\begin{tabular}{@{}l l l@{}}
\toprule
Input & Source & Updated \\
\midrule
$U_n$,\; $\vbarsla$,\; $\omega_j$ & Cluster monitor & Each round \\
$\Pmem$,\; $\Pcpu$                & Prediction layer & Each round \\
$K$,\; $\Mtax$,\; $\theta_n$,\; $\omfixed$ & Configuration & Constant \\
\bottomrule
\end{tabular}
\end{center}

% =============================================================================
\appendix
% =============================================================================

\section{How $\bar{v}_n^{\text{SLA}}$ is Computed}
\label{app:vbar}

\noindent\textbf{Additional symbols:}

\begin{center}
\begin{tabular}{@{}l p{10.5cm}@{}}
\toprule
Symbol & Meaning \\
\midrule
$K$ & Rolling window length --- number of past scheduling rounds included
  in the violation rate. Configuration scalar \\[6pt]
$k$ & Summation index, $k = 0, 1, \ldots, K{-}1$. $k=0$ is the most recent
  round; $k = K{-}1$ is the oldest \\[6pt]
$\mathbf{1}[\cdot]$ & Indicator function: 1 if the condition holds, 0
  otherwise \\
\bottomrule
\end{tabular}
\end{center}

\medskip
At scheduling round $r$:

\[
  \bar{v}_n^{\text{SLA}} \;=\;
  \frac{1}{K} \sum_{k=0}^{K-1}
    \mathbf{1}\!\left[\, U_n > \Mavail \;\text{ at round } r - k \,\right]
\]

$\vbarsla$ is the fraction of the last $K$ rounds in which actual memory on
node $n$ exceeded $\Mavail$. It is passed to the optimizer as an input each
round by the cluster monitor.

% =============================================================================
\section{How $u_n^{\text{mem}}$ is Computed}
\label{app:umem}
% =============================================================================

\[
  \umem \;=\; 1 \;+\; \min\!\left(1,\;\max\!\left(0,\;
    \frac{U_n}{\max\!\left(1,\, \Mavail\right)}\right)\right)
\]

The inner expression $U_n / \max(1, \Mavail)$ is the memory utilization
fraction --- a value in $[0,1]$ representing how full the node is. Clamping
to $[0,1]$ (via $\min(1,\,\max(0,\,\cdot))$) and adding 1 shifts the range
to $[1,2]$:

\begin{itemize}[noitemsep, topsep=4pt]
  \item $\umem = 1$: node is idle (utilization fraction $= 0$; no
    consolidation boost)
  \item $\umem = 2$: node is at full capacity (utilization fraction $= 1$;
    maximum boost)
  \item $1 < \umem < 2$: partial load (real-valued and continuous)
\end{itemize}

% =============================================================================
\section{How $\omega_j$ is Computed}
\label{app:omega}
% =============================================================================

$\omega_j$ is the priority weight for job $j$, equal to the priority weight
of the tenant that owns job $j$. The cluster monitor computes one weight per
tenant ($\omega_t$) and assigns it to each job $j$ where $t(j) = t$.

\medskip
\noindent\textbf{Additional symbols:}

\begin{center}
\begin{tabular}{@{}l p{10.5cm}@{}}
\toprule
Symbol & Meaning \\
\midrule
$\bar{W}_t$ & Average scheduling delay for tenant $t$ --- mean time from job
  submission to job assignment, over all completed jobs of tenant $t$ \\[6pt]
$\bar{W}$ & Cluster-wide average scheduling delay:
  $\displaystyle\bar{W} = \frac{1}{|T|}\sum_{t \in T} \bar{W}_t$ \\
\bottomrule
\end{tabular}
\end{center}

\[
  \omega_t \;=\; 1 \;+\; \max\!\left(0,\;
    \frac{\bar{W}_t - \bar{W}}{\max\!\left(1,\, \bar{W}\right)}
    \right)
\]

\[
  \omega_j \;=\; \omega_{t(j)}
\]

The weight is built from the excess wait above the cluster average. Adding 1
ensures the minimum is 1 (no unfairness = no extra weight). There is no stated
maximum: a tenant waiting far above average receives $\omega_j \gg 1$, giving
their jobs proportionally higher priority.

\begin{itemize}[noitemsep, topsep=4pt]
  \item $\omega_j = 1$ when $\bar{W}_{t(j)} \leq \bar{W}$ --- no boost
  \item $\omega_j > 1$ when $\bar{W}_{t(j)} > \bar{W}$ --- boost grows
    proportionally with excess wait (real-valued, no maximum)
\end{itemize}

\end{document}
